\section{Proofs}
\label{sec:appendix-proofs}

\begin{proof}[Derivation for conditional result~\ref{cond:capital-deepening-benchmark}]
Let $\widetilde k_t\equiv\widetilde K_t/L_t$. Equations
\eqref{eq:capital-benchmark-production}--
\eqref{eq:capital-benchmark-allocation} give
\begin{equation}
    \frac{\widetilde Y_t}{\widetilde K_t}
    =
    \zeta^\alpha
    \left[
        \omega_L\widetilde k_t^{-\varrho_K}
        +\omega_K(1-\zeta)^{\varrho_K}
    \right]^{(1-\alpha)/\varrho_K}.
    \label{eq:proof-capital-benchmark-output-capital}
\end{equation}
When $\varrho_K>0$, this expression is strictly decreasing in
$\widetilde k_t$, diverges as $\widetilde k_t\downarrow0$, and converges to
$B_K(\zeta)$ as $\widetilde k_t\to\infty$. Capital accumulation implies
\begin{equation}
    \frac{\dot{\widetilde k}_t}{\widetilde k_t}
    =
    s_K\frac{\widetilde Y_t}{\widetilde K_t}
    -(\delta+n).
    \label{eq:proof-capital-benchmark-law}
\end{equation}
If $s_KB_K(\zeta)<\delta+n$, monotonicity and the two boundary limits imply
a unique finite stationary value of $\widetilde k_t$. If
$s_KB_K(\zeta)>\delta+n$, equation
\eqref{eq:proof-capital-benchmark-law} converges to the positive constant
$s_KB_K(\zeta)-\delta-n$. Hence
$g_{\widetilde K}=n+g_{\widetilde k}=s_KB_K(\zeta)-\delta$; asymptotic
linearity gives the same rate for output. At equality, the right-hand side is
positive at every finite $\widetilde k_t$ and converges to zero, so
$\widetilde k_t$ diverges while its growth rate vanishes.

The inner labor contribution satisfies
$1-s^{K_2}_t\sim c\widetilde k_t^{-\varrho_K}$ for a positive constant $c$.
Since $\widetilde Y_t/L_t\sim B_K(\zeta)\widetilde k_t$, the competitive
wage obeys
\[
    w_t
    =(1-\alpha)(1-s^{K_2}_t)\frac{\widetilde Y_t}{L_t}
    \sim \widetilde c\,
    \widetilde k_t^{1-\varrho_K}
    =\widetilde c\,
    \widetilde k_t^{1/\sigma_{KL}}
\]
for a positive constant $\widetilde c$. The labor share converges to zero and,
on the positive-growth branch,
$g_w=g_{\widetilde k}/\sigma_{KL}=g_{\widetilde Y/L}/\sigma_{KL}$.
\end{proof}

\begin{proof}[Proof of Proposition~\ref{prop:wage-share-thresholds}]
The CES expenditure share in \eqref{eq:ai-composite-share} implies
\[
    \frac{\partial\log(1-s^X_t)}{\partial\log X_t}
    =
    -\left(1-\frac{1}{\sigma_{XL}}\right)s^X_t.
\]
The output elasticity follows directly from
\eqref{eq:final-production}--\eqref{eq:service-composite}. Combining these two
expressions with \eqref{eq:labor-demand} yields
\eqref{eq:labor-share-ai-elasticity} and \eqref{eq:wage-ai-elasticity}.
\end{proof}

\begin{proof}[Proof of Proposition~\ref{prop:dynamic-labor-share}]
Log-differentiating \eqref{eq:production-labor-share} and using the CES share
equation gives \eqref{eq:dynamic-production-labor-share}. Since
$L_t=(1-h_t)N_t$, $g_{L,t}=n-\dot h_t/(1-h_t)$.
Log-differentiating \eqref{eq:aggregate-labor-share} then gives
\eqref{eq:dynamic-aggregate-labor-share}.
\end{proof}

\begin{proof}[Derivation for conditional result~\ref{cond:gradual-automation}]
Equation~\eqref{eq:automation-odds} gives
\[
    \frac{s^M_t}{1-s^M_t}
    =
    \left(\frac{\omega_M}{\omega_H}\right)^{\sigma_{HM}}
    \left(\frac{w_t}{\xi}\right)^{\sigma_{HM}-1}.
\]
If $\sigma_{HM}>1$ and $w_t/\xi\to\infty$, the right-hand side diverges and
$s^M_t\to1$. Moreover, the unit cost in
\eqref{eq:research-unit-cost} converges to
\[
    P^E_t\longrightarrow
    \omega_M^{-\sigma_{HM}/(\sigma_{HM}-1)}\xi.
\]
The first expression in \eqref{eq:research-conditional-demands} then implies
$H_t/E_t\to0$. If $w_t/\xi$ remains bounded, the odds in
\eqref{eq:automation-odds} remain bounded, so $\sigma_{HM}>1$ does not by itself imply
$s^M_t\to1$.
\end{proof}

\begin{proof}[Proof of Proposition~\ref{prop:research-employment-scale}]
Take growth rates in equation~\eqref{eq:research-input-ratio} to obtain the first
identity. Combine it with equation~\eqref{eq:human-research-scale-identity} and
$g_N=n$ to obtain the second. The two sign conditions follow directly.
\end{proof}

\begin{proof}[Proof of Proposition~\ref{prop:automated-research-wage}]
Taking the $1/\sigma_{HM}$ power of
\eqref{eq:research-input-ratio} and solving for the wage gives the equality in
\eqref{eq:wage-automated-research-lower-bound}. Labor feasibility implies
$0<H_t\leq N_t$, so $M_t/H_t\geq M_t/N_t$ and gives the inequality. Because
$\xi$, the CES weights, and $\sigma_{HM}$ are positive and constant,
$M_t/N_t\to\infty$ implies $w_t\to\infty$. Log differentiation of the equality
gives \eqref{eq:wage-research-scale-growth}.

For the finite-time claim, suppose instead that $w_t$ were bounded on a terminal
interval $[t_0,T^*)$. Equation~\eqref{eq:research-input-ratio} would then bound
$M_t/H_t$. Because $H_t\leq N_t=N_0e^{nt}$ and $T^*<\infty$, both $H_t$ and $M_t$
would be bounded on that interval. The research CES would therefore bound $E_t$.
But the ideas-production equation implies
\[
    A_t^{1-\phi}
    =
    A_{t_0}^{1-\phi}
    +(1-\phi)\chi\int_{t_0}^{t}E_\tau^\eta\dd\tau.
\]
The right-hand side remains finite as $t\uparrow T^*$, contradicting
$A_t\to\infty$. Hence the wage cannot be bounded on any terminal interval, which
establishes \eqref{eq:finite-singularity-wage-unbounded}.
\end{proof}

\begin{proof}[Derivation for conditional result~\ref{cond:ai-dominated-cd-limit}]
Equations~\eqref{eq:ai-dominated-ga} and
\eqref{eq:cd-ai-growth-relations} imply
$D_{\mathrm{AI}}g_A=\eta n$. A positive finite $g_A$ therefore requires
$D_{\mathrm{AI}}>0$ and gives \eqref{eq:cd-ai-ga}--
\eqref{eq:cd-ai-gamma}. The household Euler equation, capital accumulation, the
developer's research first-order and costate conditions, and the resource constraint
then give \eqref{eq:aggregate-bg-rates}--\eqref{eq:bg-consumption-share}. The
derivation of \eqref{eq:bg-research-resource-share} is reported below. Positive
consumption is necessary for an interior household allocation.
\end{proof}

\begin{proof}[Derivation for conditional result~\ref{cond:human-essentiality-acceleration}]
Equations~\eqref{eq:fully-cd-output-capability-growth} and
\eqref{eq:cd-research-growth} imply
\eqref{eq:full-cd-ga}--\eqref{eq:full-cd-gamma}. Equations
\eqref{eq:cd-ai-ga}--\eqref{eq:cd-ai-gamma} give the AI-dominated rates. Positive
finite Cobb--Douglas rates follow from
Assumption~\ref{ass:stable-human-essential-benchmark}; positive finite
AI-dominated rates require $D_{\mathrm{AI}}>0$. Subtracting the two denominators gives
\eqref{eq:feedback-denominator-gap}. The resource constraint
gives \eqref{eq:fully-cd-consumption-feasibility} after substituting
\eqref{eq:full-cd-ga}--\eqref{eq:fully-cd-inference-share} and
\eqref{eq:full-cd-research-resource-share}.
\end{proof}

\begin{proof}[Derivation for conditional result~\ref{cond:singularity-conditions}]
Let $z_t\equiv1/g_{A,t}$. Equation
\eqref{eq:ces-transition-growth-dynamics} implies the linear equation
\[
    \dot z_t=-\eta n z_t+D(s^M_t).
\]
Multiplying by $e^{\eta nt}$ and integrating gives
\eqref{eq:time-varying-singularity-condition}. When
$\sigma_{HM}=1$, $s^M_t=\omega_M$ and the constant denominator is
$D_{\mathrm{CD}}$. If $D_{\mathrm{CD}}>0$, the logistic equation for $g_A$ has
the stable positive fixed point $\eta n/D_{\mathrm{CD}}$.

If $D_{\mathrm{CD}}=0$, direct integration gives
\[
    g_{A,t}=g_{A,0}e^{\eta nt},
    \qquad
    A_t
    =
    A_0\exp\left[
        \frac{g_{A,0}}{\eta n}
        \left(e^{\eta nt}-1\right)
    \right].
\]
Both expressions are finite for every finite $t$. If $D_{\mathrm{CD}}<0$, the
solution for the growth rate is
\[
    g_{A,t}
    =
    \frac{g_{A,0}e^{\eta nt}}
    {1-\frac{(-D_{\mathrm{CD}})g_{A,0}}{\eta n}
        \left(e^{\eta nt}-1\right)}.
\]
The denominator reaches zero at the date in
\eqref{eq:constant-feedback-singularity-time}. Near that date, $g_A$ is
proportional to $(T^*-t)^{-1}$, so both $g_A$ and $A$ diverge. When
$\sigma_{HM}>1$, $D(s^M_t)$ is time varying. If it remains nonnegative, the bracket
in \eqref{eq:time-varying-singularity-condition} remains positive. Otherwise the
same divergence occurs exactly when that bracket reaches zero.
\end{proof}

\begin{proof}[Derivation for conditional result~\ref{cond:production-ces-regimes}]
Let $x_t\equiv X_t/L_t$, $k_t\equiv K_t/L_t$, and write
$s_t\equiv s^X_t$. Final output per production worker is
\[
    \frac{Y_t}{L_t}
    =
    k_t^\alpha
    \left[
        \omega_L
        +\omega_Xx_t^{(\sigma_{XL}-1)/\sigma_{XL}}
    \right]^{(1-\alpha)\sigma_{XL}/(\sigma_{XL}-1)}.
\]
Using this expression in the monopoly first-order condition gives
\begin{equation}
    \frac{\xi}{A_tk_t^\alpha}
    =
    (1-\alpha)s_t[1-\varepsilon(s_t)]
    \frac{
        \left[
            \omega_L
            +\omega_Xx_t^{(\sigma_{XL}-1)/\sigma_{XL}}
        \right]^{(1-\alpha)\sigma_{XL}/(\sigma_{XL}-1)}
    }{x_t},
    \qquad
    \varepsilon(s)
    =
    \frac{1-s}{\sigma_{XL}}+\alpha s.
    \label{eq:appendix-normalized-monopoly-condition}
\end{equation}
Condition~\eqref{eq:vanishing-normalized-ai-cost} makes the left-hand side
converge to zero.

For $\sigma_{XL}<1$, revenue approaches zero both as $x_t\to0$ and as
$x_t\to\infty$. As normalized marginal cost converges to zero, the monopolist
approaches the interior revenue-maximizing quantity, which satisfies
$\varepsilon(\bar s^X)=1$.
Solving this condition gives
\[
    \bar s^X=\frac{1-\sigma_{XL}}{1-\alpha\sigma_{XL}}.
\]
Because the CES mapping between $s^X$ and $X/L$ is one-to-one, the corresponding
service ratio is finite.

For $\sigma_{XL}=1$, the CES contribution $s^X_t$ equals the constant weight $\omega_X$.

For $\sigma_{XL}>1$, the right-hand side of
\eqref{eq:appendix-normalized-monopoly-condition} is strictly positive at every
finite $x_t>0$ and diverges as $x_t\to0$. It can therefore converge to zero only if
$x_t\to\infty$. The CES share equation then implies $s^X_t\to1$. Consequently,
$\varepsilon^X_t\to\alpha$. Multiplying the monopoly first-order condition by
$X_t/Y_t$ and using $X_t=A_tU_t$ gives
\[
    \frac{\xi U_t}{Y_t}\longrightarrow(1-\alpha)^2.
\]
Since
\[
    Z_t\sim
    \omega_X^{\sigma_{XL}/(\sigma_{XL}-1)}X_t,
\]
substituting $X_t=A_tU_t$ and the limiting compute share into production yields
\[
    \frac{Y_t}{K_t}
    \sim
    \left[
        \frac{\omega_X^{\sigma_{XL}/(\sigma_{XL}-1)}(1-\alpha)^2}{\xi}
    \right]^{(1-\alpha)/\alpha}
    A_t^{(1-\alpha)/\alpha}.
\]
This establishes \eqref{eq:high-sigma-ak-limit}.
\end{proof}

\begin{proof}[Proof of Proposition~\ref{prop:production-unit-elasticity-knife-edge}]
Combining relative demand in \eqref{eq:relative-ai-demand} with the definition of
$s^X_t$ gives
\[
    \frac{s^X_t}{1-s^X_t}
    =
    \frac{\omega_X}{\omega_L}
    \left(\frac{X_t}{L_t}\right)^{(\sigma_{XL}-1)/\sigma_{XL}}
    =
    \left(\frac{\omega_X}{\omega_L}\right)^{\sigma_{XL}}
    \left(\frac{w_t}{p^X_t}\right)^{\sigma_{XL}-1}.
\]
For every finite relative price $\mathcal R>0$, the right-hand side converges to
$\omega_X/\omega_L$ as $\sigma\to1$, which implies
$s^X(\mathcal R,\sigma)\to\omega_X$. At $\sigma=1$, the exponent on the relative
price is zero and the share is identically $\omega_X$. For every fixed $\sigma>1$,
the odds diverge as $\mathcal R\to\infty$, so
$s^X(\mathcal R,\sigma)\to1$. These observations establish the two
iterated limits in \eqref{eq:noncommuting-production-limits}.

Finally, set $s^X=\bar s$ in \eqref{eq:production-share-odds} and solve for
$\mathcal R$ to obtain \eqref{eq:production-share-threshold}. Because
$\bar s>\omega_X$,
\[
    \log\frac{\bar s}{1-\bar s}
    +\log\frac{\omega_L}{\omega_X}>0.
\]
The logarithm of \eqref{eq:production-share-threshold} therefore diverges at rate
$1/(\sigma-1)$ as $\sigma\downarrow1$.
\end{proof}

\begin{proof}[Proof of Proposition~\ref{prop:gross-substitution-no-steady-state}]
Suppose, to obtain a contradiction, that such a steady state exists and that
$g_{A,t}\to\bar g_A>0$. Then $A_t\to\infty$.
The CES share definition gives the exact identity
\begin{equation}
    \frac{Z_t}{X_t}
    =
    \left(\frac{\omega_X}{s^X_t}\right)^{
        \sigma_{XL}/(\sigma_{XL}-1)}.
    \label{eq:appendix-z-x-share-identity}
\end{equation}
The monopoly first-order condition and final-good demand for AI services imply
\begin{equation}
    \frac{X_t}{Y_t}
    =
    \frac{A_t}{\xi}(1-\alpha)s^X_t
    (1-\varepsilon^X_t).
    \label{eq:appendix-x-y-monopoly-identity}
\end{equation}
Since
\[
    \frac{Y_t}{K_t}
    =
    \left(
        \frac{Z_t}{X_t}
        \frac{X_t}{Y_t}
        \frac{Y_t}{K_t}
    \right)^{1-\alpha},
\]
substitution of \eqref{eq:appendix-z-x-share-identity} and
\eqref{eq:appendix-x-y-monopoly-identity}, followed by solving for $Y_t/K_t$,
yields
\begin{equation}
    \frac{Y_t}{K_t}
    =
    \left[
        \frac{1-\alpha}{\xi}
        \omega_X^{\sigma_{XL}/(\sigma_{XL}-1)}
        (s^X_t)^{-1/(\sigma_{XL}-1)}
        (1-\varepsilon^X_t)A_t
    \right]^{(1-\alpha)/\alpha}.
    \label{eq:exact-high-sigma-output-capital}
\end{equation}
This equality makes clear which parts of the result require AI dominance. If
$s^X_t\to1$, then $\varepsilon^X_t\to\alpha$, and
\eqref{eq:exact-high-sigma-output-capital} becomes
\begin{equation}
    \frac{Y_t}{K_t}
    \sim
    D_{XL}A_t^{(1-\alpha)/\alpha},
    \qquad
    D_{XL}
    =
    \left[
        \frac{(1-\alpha)^2}{\xi}
        \omega_X^{\sigma_{XL}/(\sigma_{XL}-1)}
    \right]^{(1-\alpha)/\alpha}>0.
    \label{eq:appendix-high-sigma-output-capital-limit}
\end{equation}
Conditional result~\ref{cond:production-ces-regimes} establishes $s^X_t\to1$ under
\eqref{eq:vanishing-normalized-ai-cost}.

The nonexistence result does not require that additional condition. When
$\sigma_{XL}>1$ and $s^X_t\in(0,1)$,
\[
    (s^X_t)^{-1/(\sigma_{XL}-1)}\geq1,
    \qquad
    1-\varepsilon^X_t
    \geq
    \min\left\{1-\alpha,1-\frac{1}{\sigma_{XL}}\right\}>0.
\]
Equation~\eqref{eq:exact-high-sigma-output-capital} therefore implies
\begin{equation}
    \frac{Y_t}{K_t}
    \geq
    \underline D_{XL}A_t^{(1-\alpha)/\alpha},
    \qquad
    \underline D_{XL}>0.
    \label{eq:high-sigma-output-capital-lower-bound}
\end{equation}
If $A_t\to\infty$, capital demand gives
\begin{equation}
    r_t
    =\alpha\frac{Y_t}{K_t}
    \longrightarrow\infty.
    \label{eq:high-sigma-return-divergence}
\end{equation}
Because $c_t=C_t/N_t$ and $g_N=n$, the Euler equation implies
\begin{equation}
    g_{C,t}
    =n+r_t-\delta-\rho
    \longrightarrow\infty.
    \label{eq:high-sigma-consumption-growth-divergence}
\end{equation}
Aggregate consumption cannot consequently grow at a finite constant rate. If $g_A$
converged to the positive constant $\bar g_A$, the alleged steady state would require
both finite constant consumption growth and divergent consumption growth. This is the
required contradiction.
\end{proof}

\begin{proof}[Derivation for conditional result~\ref{cond:high-sigma-singular-scaling}]
Let $z_t\equiv Y_t/K_t$, $u_t\equiv q_tA_t/K_t$, and
$a_t\equiv g_{A,t}/z_t$. AI dominance and
\eqref{eq:appendix-high-sigma-output-capital-limit} imply
\begin{equation}
    z_t\sim D_{XL}A_t^\vartheta,
    \qquad
    \frac{\xi U_t}{Y_t}\longrightarrow d_U,
    \qquad
    \frac{r_t}{z_t}\longrightarrow\alpha.
    \label{eq:appendix-singular-production-limits}
\end{equation}
Moreover,
\begin{equation}
    \frac{\pi^X_{A,t}/q_t}{z_t}
    =
    \frac{A_t\pi^X_{A,t}}{Y_t}\frac{1}{u_t}
    \longrightarrow
    \frac{d_U}{u_t}.
    \label{eq:appendix-singular-profit-value}
\end{equation}

Let $\bar i$, $\bar c$, $\bar m_R$, $\bar a$, and $\bar u$ denote the
positive limits of $I/Y$, $C/Y$, $\xi M/Y$, $a_t$, and $u_t$. The household
Euler equation gives $g_C/z\to\alpha$. Asymptotic regularity of $C/Y$ therefore
gives $g_Y/z\to\alpha$. Because
$g_z=\vartheta g_A$ and $g_K/z\to\bar i$, output growth also satisfies
\[
    \frac{g_Y}{z}
    \longrightarrow
    \bar i+\vartheta\bar a.
\]
It follows that
\begin{equation}
    \bar i=\alpha-\vartheta\bar a.
    \label{eq:appendix-singular-investment}
\end{equation}

When $s^M_t\to1$, the automated-research first-order condition becomes
$\xi=q_t\eta\dot A_t/M_t$ asymptotically. Hence
\begin{equation}
    \bar m_R
    =\eta\bar u\bar a.
    \label{eq:appendix-singular-research-share}
\end{equation}
The growth rate of $u_t=q_tA_t/K_t$ and the developer's costate equation give
\begin{align}
    \frac{g_{u,t}}{z_t}
    &\longrightarrow
    \alpha-\bar i+(1-\phi)\bar a-\frac{d_U}{\bar u}\notag\\
    &=
    (1-\phi+\vartheta)\bar a-\frac{d_U}{\bar u}.
    \label{eq:appendix-singular-shadow-ratio}
\end{align}
Asymptotic regularity of $u_t$ sets the left-hand side of
\eqref{eq:appendix-singular-shadow-ratio} equal to zero, so
\begin{equation}
    \bar u
    =
    \frac{d_U}{(1-\phi+\vartheta)\bar a}.
    \label{eq:appendix-singular-shadow-limit}
\end{equation}

Finally, automated research dominance and asymptotic regularity of the
research-resource share imply
$g_M/z\to g_Y/z\to\alpha$. Log differentiation of the ideas-production function
therefore gives
\[
    \frac{g_{g_A}}{z}
    \longrightarrow
    (\phi-1)\bar a+\eta\alpha.
\]
Asymptotic regularity of $a_t=g_A/z$, together with
$g_z=\vartheta g_A$, implies that the left-hand side also converges to
$\vartheta\bar a$. Equating the two expressions yields
\[
    \bar a
    =
    \frac{\eta\alpha}{1-\phi+\vartheta}.
\]
Equations~\eqref{eq:appendix-singular-research-share} and
\eqref{eq:appendix-singular-shadow-limit} then give
$\bar u=d_U/(\eta\alpha)$ and
$\bar m_R=\eta d_U/(1-\phi+\vartheta)$. The resource constraint gives
$\bar c=1-d_U-\bar m_R-\bar i$.

To obtain the time path, use $g_z=\vartheta g_A\sim\vartheta\bar a z$:
\[
    \dot z_t
    \sim
    \vartheta\bar a z_t^2.
\]
Integration gives
$z_t\sim[\vartheta\bar a(T^*-t)]^{-1}$. Finally,
\[
    \frac{K_t}{C_t}
    =
    \frac{1}{(C_t/Y_t)z_t}
    \longrightarrow0,
    \qquad
    \frac{q_tA_t}{C_t}
    =
    u_t\frac{K_t}{C_t}
    \longrightarrow0.
\]
These limits establish the terminal-value claims.

For the wage claim, combine $g_Y/z\to\alpha$ with the asymptotic expression for
$z_t$. An application of l'H\^opital's rule gives
\[
    \frac{\log Y_t}{\log[1/(T^*-t)]}
    \longrightarrow
    \frac{\alpha}{\vartheta\bar a}>0.
\]
Thus $Y_t\to\infty$. Population remains finite because
$N_t=N_0e^{nt}$ and $T^*<\infty$. Equation~\eqref{eq:appendix-singular-research-share}
and constant $\xi$ imply $M_t/Y_t\to\bar m_R/\xi>0$. Hence
$Y_t/N_t\to\infty$ and $M_t/N_t\to\infty$. Proposition
\ref{prop:automated-research-wage} then implies $w_t\to\infty$.
\end{proof}

\begin{proof}[Derivation for conditional result~\ref{cond:nonunit-ces-steady-states}]
Consider first $\sigma_{XL}<1$. Conditional result~\ref{cond:production-ces-regimes}
implies that $X_t/L_t$ converges to a positive constant. Because $L/N$ also
converges to a positive constant, $g_L=g_X=n$.

Because $K/Y$ is constant, $g_K=g_Y$. Because $X/L$ is constant, the service
composite grows at rate $n$. The final-good production function therefore gives
\[
    g_Y=\alpha g_Y+(1-\alpha)n,
\]
so $g_Y=g_K=n$. Constant $C/Y$ gives $g_C=n$, and the household Euler equation
implies $r-\delta=\rho$. Labor demand then implies $g_w=0$. Finally, $X=AU$ gives
$g_U=n-g_A$.

Let $m$ denote the common asymptotic growth rate of the two research inputs. This
common rate follows from \eqref{eq:research-input-ratio}: the real wage and $\xi$ are
constant, so $H_t/M_t$ is constant. Constant returns in the research CES then give
$g_E=m$, while labor feasibility implies $m\leq n$.

Write $a\equiv g_A>0$. Constant capability growth and $g_E=m$ imply
\begin{equation}
    (1-\phi)a=\eta m.
    \label{eq:appendix-low-sigma-ideas-growth}
\end{equation}
The research first-order condition $q_tF_{M,t}=\xi$ supplies a second restriction.
Because $g_{F_M}=a-m$, it implies
\begin{equation}
    g_q=m-a.
    \label{eq:appendix-low-sigma-shadow-value-growth}
\end{equation}
At the optimal service quantity, the envelope theorem gives
\begin{equation}
    \pi^X_{A,t}=\frac{\xi X_t}{A_t^2},
    \qquad
    g_{\pi^X_A}=n-2a.
    \label{eq:appendix-low-sigma-capability-profit}
\end{equation}
Divide the developer's costate equation \eqref{eq:private-costate} by $q_t$. Since
$F_A=\phi a$, equations~\eqref{eq:appendix-low-sigma-ideas-growth} and
\eqref{eq:appendix-low-sigma-shadow-value-growth} give
\begin{equation}
    \rho
    =
    \frac{\pi^X_{A,t}}{q_t}
    +\phi a+m-a
    =
    \frac{\pi^X_{A,t}}{q_t}+(1-\eta)m.
    \label{eq:appendix-low-sigma-costate-limit}
\end{equation}
Because $\rho>n$ and $m\leq n$, the ratio $\pi^X_A/q$ converges to the strictly
positive constant $\rho-(1-\eta)m$. Its growth rate must therefore converge to zero.
Equations~\eqref{eq:appendix-low-sigma-shadow-value-growth} and
\eqref{eq:appendix-low-sigma-capability-profit} imply
\[
    g_{\pi^X_A/q}=n-2a-(m-a)=n-m-a,
\]
so $m+a=n$. Combining this condition with
\eqref{eq:appendix-low-sigma-ideas-growth} yields
\[
    a=\frac{\eta n}{1-\phi+\eta},
    \qquad
    m=\frac{(1-\phi)n}{1-\phi+\eta}.
\]
It follows that $g_{H/N}=m-n=-a$ and, because $g_M=m$, $g_Y=n$, and $\xi$ is
constant, $g_{\xi M/Y}=-a$. This proves part (i). Part (ii) is
the exact lower-bound argument in the proof of
Proposition~\ref{prop:gross-substitution-no-steady-state}, with unbounded capability
imposed as part of the conditional candidate.
\end{proof}

\begin{proof}[Derivation of Equation~\eqref{eq:bg-research-resource-share}]

In the Cobb--Douglas production benchmark, the envelope theorem and
\eqref{eq:cd-monopoly} imply
\[
    \pi^X_{A,t}=\frac{\beta^2Y_t}{A_t}.
\]
The research first-order condition in the AI-dominated limit gives
\[
    \xi=q_t\eta\frac{\dot A_t}{M_t},
    \qquad
    \frac{q_tA_t}{Y_t}
    =
    \frac{m_R^\infty}{\eta g_A^\infty}.
\]
Because the right-hand side is constant on a balanced-growth path,
$g_q=g_Y-g_A$. Dividing \eqref{eq:private-costate} by $q_t$, using
$r-\delta=\rho+\gamma^\infty$ and $g_Y=n+\gamma^\infty$, and substituting the
two expressions above gives
\[
    \rho-n
    =
    \frac{\beta^2\eta g_A^\infty}{m_R^\infty}
    -(1-\phi)g_A^\infty.
\]
Solving for $m_R^\infty$ yields
\eqref{eq:bg-research-resource-share}.
\end{proof}

\begin{proof}[Derivation for conditional result~\ref{cond:sigma-hm-one-high-sigma-xl}]
The first claim follows directly from
\eqref{eq:sigma-hm-one-high-sigma-xl-capital-growth}. For the second, set
$x_t=\log A_t$, $k_t=\log K_t$, $b=\eta\omega_M$, and
$d=1-\phi-\vartheta b$. After using \eqref{eq:high-sigma-ak-limit} and dropping terms that
vanish in the high-capability limit, the reduced system can be written as
\begin{align}
    \dot x_t
    &=C\exp\!\left\{\eta\omega_Hnt+bk_t-dx_t\right\},
    \label{eq:sigma-hm-one-high-sigma-xl-x}\\
    \dot k_t
    &=S\exp(\vartheta x_t)-\delta,
    \qquad C,S>0.
    \label{eq:sigma-hm-one-high-sigma-xl-k}
\end{align}
Define $u_t\equiv\dot x_t\exp(-\vartheta x_t)$. In this asymptotic system,
\begin{equation}
    \frac{\dot u_t}{u_t}
    =
    \eta\omega_Hn-b\delta
    +\left[bS-(d+\vartheta)u_t\right]\exp(\vartheta x_t).
    \label{eq:sigma-hm-one-high-sigma-xl-u}
\end{equation}
By the maintained assumption $A_t\to\infty$, $x_t$ eventually becomes arbitrarily
large. Because $b>0$, equation~\eqref{eq:sigma-hm-one-high-sigma-xl-u} then keeps $u_t$
bounded away from zero, whether $d+\vartheta$ is positive, zero, or negative. It follows that
\begin{equation}
    \frac{\dd}{\dd t}\exp(-\vartheta x_t)=-\vartheta u_t
\end{equation}
is eventually bounded above by a negative constant. Hence $\exp(-\vartheta x_t)$ reaches
zero in finite time, so $A_t$ diverges at a finite date.
\end{proof}

\begin{proof}[Derivation for conditional result~\ref{cond:knowledge-wedge}]
Evaluate the planner's marginal net benefit from research compute at the private
allocation. Using \eqref{eq:private-compute-foc},
\[
    Q_tF_{M,t}-\xi
    =
    (Q_t-q_t)F_{M,t}.
\]
If $Q_t>q_t$, the expression is positive and the planner increases research compute
locally relative to the private allocation.
\end{proof}
